\lysh{13176}{4}

\begin{alist}
\item Έχουμε: $|x - 1| < 2 \Leftrightarrow -2 < x - 1 < 2 \Leftrightarrow -1 < x < 3$.

Για να λύσουμε την ανίσωση $x^2 - 3x + 2 \ge 0$, θα βρούμε πρώτα τις ρίζες $x_1$ και $x_2$ του τριωνύμου. Έχουμε λοιπόν: $x_1 + x_2 = -\dfrac{\beta}{\alpha} = 3$ και $x_1 \cdot x_2 = \dfrac{\gamma}{\alpha} = 2$, οπότε $x_1 = 1$ και $x_2 = 2$.
Στη συνέχεια κατασκευάζουμε τον πίνακα προσήμου για το τριώνυμο $x^2 - 3x + 2$, με $\alpha = 1 > 0$:
\begin{center}
\begin{tabular}{|c|ccccccc|}

$x$ & $-\infty$ & & $1$ & & $2$ & & $+\infty$ \\ 
$x^2 - 3x + 2$ & & $+$ & $0$ & $-$ & $0$ & $+$ & \\ 
\end{tabular}
\end{center}
Οπότε η ανίσωση αληθεύει για $x \in (-\infty, 1] \cup [2, +\infty)$.

\item Με χρήση του άξονα των πραγματικών αριθμών βλέπουμε ότι οι ανισώσεις συναληθεύουν για
$x \in (-1, 1] \cup [2, 3)$.

\item Εφόσον οι αριθμοί $\rho_1$ και $\rho_2$, με $\rho_1 < \rho_2$, ανήκουν στο σύνολο των κοινών λύσεων των ανισώσεων, θα ισχύει: $\rho_1, \rho_2 \in (-1, 1] \cup [2, 3)$.
\begin{rlist}
\item Αν $\rho_1, \rho_2 \in (-1, 1]$, τότε: $\begin{cases} -1 < \rho_1 \le 1 \\ -3 < 3\rho_2 \le 3 \end{cases}$ και προσθέτοντας κατά μέλη προκύπτει:
$-4 < \rho_1 + 3\rho_2 \le 4$, συνεπώς $-1 < \dfrac{\rho_1 + 3\rho_2}{4} \le 1$ και ο αριθμός είναι κοινή λύση των ανισώσεων.
\item Αν $\rho_1 \in (-1, 1]$ και $\rho_2 \in [2, 3)$, τότε $\begin{cases} -1 < \rho_1 \le 1 \\ 6 \le 3\rho_2 < 9 \end{cases}$, και προσθέτοντας κατά μέλη προκύπτει: $5 < \rho_1 + 3\rho_2 < 10$, συνεπώς $\dfrac{5}{4} < \dfrac{\rho_1 + 3\rho_2}{4} < \dfrac{5}{2}$ και ο αριθμός είναι κοινή λύση των ανισώσεων μόνο εάν $2 \le \dfrac{\rho_1 + 3\rho_2}{4} < \dfrac{5}{2}$.
\end{rlist}
\end{alist}
